\section{Kth Smallest Element in a BST}
\label{sec:bst-kth-smallest}

\problemstatement{Given the root of a BST and an integer $k$, find the
$k$-th smallest value stored in the tree.}

\begin{patternbox}
\emph{``Kth smallest''} on a BST is answered directly by
Section~\ref{sec:bst-intro}'s headline fact: an \textbf{inorder
traversal visits values in strictly increasing order}. So the $k$-th
value produced by an inorder traversal \emph{is} the $k$-th smallest
value in the tree -- no sorting, and no need to visit the whole tree if
we stop as soon as the $k$-th value is produced.
\end{patternbox}

\subsection*{Understanding the problem carefully}

A recursive inorder traversal could be used, but stopping it early (the
moment the $k$-th value is produced) is awkward with plain recursion,
since recursive calls don't have a natural way to ``stop everything''
once we are several levels deep. An \textbf{iterative} inorder traversal,
using an explicit stack to simulate the recursion, gives full control over
exactly when to stop -- we simply count nodes as we pop them and halt the
moment the count reaches $k$.

\begin{ideabox}[Iterative Inorder Traversal, Counting as We Go]
Maintain an explicit stack. Repeatedly: push the current node and move
left, until reaching null (this pushes the entire leftward chain, since
inorder visits left subtrees first). Then pop a node -- this is the next
value in sorted order; increment a counter, and if the counter reaches
$k$, return this node's value immediately. Otherwise, move to the popped
node's right child and repeat the whole process (push its leftward
chain, then continue popping).
\end{ideabox}

\subsection*{Why the iterative version visits nodes in exactly inorder sequence}

Pushing the entire leftward chain before popping anything ensures that,
when we do pop, we are popping the smallest value not yet visited: any
node still left to visit that is smaller than the current pop candidate
would have to be in some left subtree along the path already pushed, and
all of those are already on the stack, ready to be popped first (stacks
pop most-recently-pushed first, which corresponds to the deepest,
leftmost, smallest unvisited node). After popping a node, moving to its
right child and repeating correctly resumes the traversal exactly where
recursive inorder would have continued, since everything in that node's
own left subtree has already been fully visited.

\subsection*{Algorithm}

\begin{enumerate}
  \item Initialize an empty stack, $\textit{count} \gets 0$,
        $\textit{node} \gets \textit{root}$.
  \item While $\textit{node}$ is not null or the stack is not empty:
  \begin{enumerate}
    \item While $\textit{node}$ is not null: push it; $\textit{node}
          \gets \textit{node}.\textit{left}$.
    \item Pop $\textit{node}$; $\textit{count} \gets \textit{count}+1$.
    \item If $\textit{count} = k$: return $\textit{node}.\textit{val}$.
    \item $\textit{node} \gets \textit{node}.\textit{right}$.
  \end{enumerate}
\end{enumerate}

\subsection*{Dry run}

\dryrun{Take the BST with root $5$, left child $3$ (children $1,4$),
right child $8$, $k=3$.}

\begin{itemize}
  \item Push leftward chain from $5$: push $5$, push $3$, push $1$ (no
        left child). Stack (top to bottom) $=[1,3,5]$.
  \item Pop $1$: count$=1 \ne 3$. Move to $1$'s right (null).
  \item Push leftward chain from null: nothing to push.
  \item Pop $3$: count$=2 \ne 3$. Move to $3$'s right, which is $4$.
  \item Push leftward chain from $4$: push $4$ (no left child).
  \item Pop $4$: count$=3=k$. Return $4$.
\end{itemize}

The $3$rd smallest value is $4$ (inorder sequence: $1,3,4,5,8$).

\subsection*{C++ implementation}

\begin{lstlisting}
#include <bits/stdc++.h>
using namespace std;

struct TreeNode {
    int val;
    TreeNode* left;
    TreeNode* right;
    TreeNode(int v) : val(v), left(nullptr), right(nullptr) {}
};

int kthSmallest(TreeNode* root, int k) {
    stack<TreeNode*> st;
    TreeNode* node = root;
    int count = 0;

    while (node || !st.empty()) {
        while (node) {
            st.push(node);
            node = node->left;
        }
        node = st.top();
        st.pop();
        count++;
        if (count == k) return node->val;
        node = node->right;
    }
    return -1; // k larger than tree size
}

int main() {
    TreeNode* root = new TreeNode(5);
    root->left = new TreeNode(3);
    root->right = new TreeNode(8);
    root->left->left = new TreeNode(1);
    root->left->right = new TreeNode(4);

    cout << "3rd smallest: " << kthSmallest(root, 3) << endl;
    return 0;
}
\end{lstlisting}

\begin{complexitybox}
\textbf{Time Complexity.} $O(h+k)$: pushing the initial leftward chain
costs $O(h)$, and each subsequent step processes one more node until the
$k$-th is found, giving $O(h+k)$ overall -- in the worst case $O(n)$ if
$k$ is close to the total node count.
\textbf{Space Complexity.} $O(h)$ for the stack.
\end{complexitybox}

\begin{takeawaybox}
Whenever a BST problem needs the $i$-th value in sorted order, or values
processed one at a time in sorted order with the ability to stop early,
an \emph{iterative} inorder traversal with an explicit stack is the right
tool -- it gives the same sorted-order guarantee as recursive inorder,
with full control over when to stop. This exact iterative inorder
machinery, generalized to resume across multiple separate calls, is what
Section~\ref{sec:bst-iterator}'s BST Iterator formalizes next.
\end{takeawaybox}
